\chapter{Kesimpulan dan Saran}
\label{chap:kesimpulan}

\section{Kesimpulan}
\label{sec:kesimpulan}

Berdasarkan hasil implementasi, pengujian, dan analisis yang telah dilakukan terhadap dua pendekatan penyelesaian permasalahan \textit{Colored Queens} pada tugas akhir ini, dapat ditarik beberapa kesimpulan sebagai berikut:

\begin{enumerate}
	\item Perangkat lunak permainan \textit{Colored Queens} berhasil dibangun menggunakan bahasa pemrograman Java dengan antarmuka grafis berbasis Java Swing. Perangkat lunak ini terdiri atas tiga halaman utama, yaitu halaman beranda, halaman pemilihan level, dan halaman permainan, yang menyajikan papan permainan dengan ukuran beragam mulai dari $7 \times 7$ hingga $30 \times 30$ yang dimuat dari berkas berformat JSON.
	
	\item Solusi permainan \textit{Colored Queens} berhasil dibangun menggunakan dua pendekatan algoritmik yang berbeda. Pendekatan pertama adalah \textit{backtracking} sistematis yang diperkuat dengan \textit{Forward Checking} dan propagasi kendala AC-3, yang merepresentasikan permasalahan sebagai \textit{Constraint Satisfaction Problem} (CSP) dengan variabel berupa sektor warna, domain berupa sel-sel pada sektor tersebut, dan kendala yang mencakup larangan baris, kolom, dan ketetanggaan dalam delapan arah. Pendekatan kedua adalah \textit{Particle Swarm Optimization} (PSO) diskrit dengan representasi indeks sel per warna, fungsi \textit{fitness} berbasis pelanggaran kendala, dan topologi \textit{neighborhood} lokal (lBest) yang menerapkan adaptasi formula kecepatan dan posisi PSO untuk ruang pencarian diskrit.
	
	\item Kedua solver berhasil diintegrasikan ke dalam perangkat lunak permainan sebagai kelas independen yang menerima objek papan sebagai masukan dan menghasilkan penugasan menteri sebagai keluaran. Integrasi diwujudkan melalui mekanisme eksekusi asinkron, sehingga solver dapat dijalankan di latar belakang untuk mendukung fitur \textit{hint} dan \textit{auto-solve} tanpa menghambat responsivitas antarmuka pengguna.
	
	\item Berdasarkan hasil pengujian, kedua solver menunjukkan kinerja yang sangat berbeda dalam mencari solusi permainan \textit{Colored Queens}. Algoritma \textit{backtracking} dengan \textit{Forward Checking} dan AC-3 mampu menyelesaikan seluruh instansi yang diujikan, mulai dari ukuran $7 \times 7$ hingga $30 \times 30$, dengan tingkat keberhasilan 100\% dan waktu penyelesaian yang cepat dan deterministik. Sebaliknya, algoritma PSO diskrit hanya mampu menemukan solusi yang valid pada papan berukuran kecil ($7 \times 7$ hingga $9 \times 9$) dengan konfigurasi parameter tertentu, namun mengalami stagnasi dan kegagalan konvergensi yang signifikan pada papan berukuran $10 \times 10$ ke atas.
	
	\item Hasil eksperimen mengonfirmasi bahwa kombinasi mekanisme propagasi kendala (\textit{Forward Checking} dan AC-3) secara signifikan mempersempit ruang pencarian melalui eliminasi nilai-nilai domain yang tidak konsisten sebelum penelusuran dilanjutkan. Pada PSO, eksperimen variasi parameter menunjukkan bahwa peningkatan jumlah partikel dan pengurangan jumlah \textit{neighborhood} memberikan dampak positif terbesar pada kemampuan eksplorasi, namun bahkan kombinasi parameter optimal tidak mampu menjamin keberhasilan pada instansi berukuran besar. Dengan demikian, untuk permasalahan \textit{Colored Queens} secara khusus dan permasalahan CSP dengan kendala keras secara umum, pendekatan berbasis propagasi kendala terbukti jauh lebih efektif dibandingkan pendekatan \textit{metaheuristik} berbasis populasi.
\end{enumerate}

\section{Saran}
\label{sec:saran}

Berdasarkan temuan dan keterbatasan yang diidentifikasi dalam penelitian ini, beberapa saran untuk pengembangan selanjutnya dapat dikemukakan.

Pertama, penyetelan parameter PSO dapat dilakukan secara lebih sistematis dan komprehensif. Eksperimen pada tugas akhir ini hanya mengeksplorasi 13 konfigurasi parameter dengan mengubah satu parameter pada satu waktu (\textit{one-factor-at-a-time}), sehingga interaksi antarparameter belum tertangkap secara penuh. Penelitian selanjutnya dapat menerapkan metode penyetelan yang lebih sistematis seperti \textit{grid search}, \textit{random search}, atau \textit{Bayesian optimization}, guna mengidentifikasi kombinasi parameter yang benar-benar optimal dan mengevaluasi potensi PSO pada permasalahan \textit{Colored Queens} secara lebih adil.

Kedua, mekanisme rekonfigurasi \textit{neighborhood} secara adaptif dapat ditambahkan ke dalam solver PSO untuk menangani stagnasi. Pada implementasi saat ini, struktur \textit{neighborhood} bersifat statis sepanjang proses pencarian. Penelitian selanjutnya dapat mengeksplorasi mekanisme yang secara dinamis mengonfigurasi ulang \textit{neighborhood} saat stagnasi terdeteksi, sehingga partikel memiliki kesempatan untuk keluar dari \textit{local optima} dan mengeksplorasi area pencarian yang berbeda.

Ketiga, pendekatan \textit{hybrid} yang menggabungkan propagasi kendala dengan metaheuristik dapat dieksplorasi sebagai arah penelitian lebih lanjut. Sebagai contoh, PSO dapat digunakan untuk menghasilkan solusi awal yang kemudian diperbaiki oleh \textit{backtracking}, atau propagasi kendala dapat diterapkan sebagai mekanisme perbaikan \textit{fitness} di dalam PSO. Pendekatan semacam ini berpotensi mengatasi kelemahan konvergensi PSO pada permasalahan dengan kendala keras.

Keempat, pengembangan heuristik pemilihan variabel yang lebih adaptif pada \textit{backtracking} dapat menjadi arah penelitian yang menjanjikan. Implementasi \textit{backtracking} dalam penelitian ini menggunakan pengurutan warna berdasarkan ukuran domain secara statik. Heuristik dinamis seperti \textit{Minimum Remaining Values} (MRV) atau \textit{Least Constraining Value} (LCV) yang memperbarui urutan pemilihan variabel secara adaptif selama proses pencarian berpotensi mengurangi jumlah \textit{backtrack} lebih lanjut.

Kelima, mengingat keterbatasan PSO yang ditemukan, penelitian selanjutnya dapat membandingkan PSO dengan metaheuristik lain seperti \textit{Simulated Annealing}, \textit{Genetic Algorithm}, atau \textit{Tabu Search}, agar dapat menilai apakah kelemahan yang ditemukan bersifat umum pada seluruh pendekatan \textit{metaheuristik}, atau merupakan karakteristik spesifik PSO pada permasalahan ini.

Keenam, pengujian pada instansi yang lebih beragam perlu dilakukan. Penelitian ini dibatasi hingga ukuran $30 \times 30$ dengan satu instansi untuk ukuran besar, sehingga penelitian selanjutnya sebaiknya mencakup lebih banyak instansi dengan variasi tata letak warna yang berbeda guna memperoleh gambaran statistik yang lebih komprehensif.